\chapter{Existing Approaches for CNOT Gate Composition}
\label{ch:existing-approaches}

In this section, we review why CNOT gate composition can be regarded as an
elementary matrix operation, and then discuss existing methods for CNOT
circuit synthesis that have aimed to optimize circuit depth.

\section{CNOT Circuits and Elementary Matrix Operations}
\label{sec:CNOT-transformation-to-Matrix-Operations}
Before reviewing existing synthesis methods, we first establish the mathematical correspondence between CNOT circuits and matrix operations over $\mathbb{F}_{2}$, which underlies all of the approaches discussed in this chapter.

The foundational result connecting these operations to CNOT circuit
synthesis is given by Theorem~\ref{thm:basic-elementary-matrix-operation-type-1-and-3}. To transform the
matrix $A_{n \times n}$ into a permutation matrix $P_{n \times n}$, we
consider elementary row and column operations. Applying a sequence of
elementary row operations $R_{1}, \ldots, R_{k}$ to $A$ yields a
decomposition of $A$; since every quantum gate is invertible, the inverses
of these elementary row operations, $R_{1}^{-1}, \ldots, R_{k}^{-1}$, can
likewise be applied to $P$ to recover $A$. Similarly, elementary column
operations $C_{1}, \ldots, C_{k'}$ can be applied to $A$, and their inverses
$C_{1}^{-1}, \ldots, C_{k'}^{-1}$ exist as well.

\begin{theorem}
\label{thm:basic-elementary-matrix-operation-type-1-and-3}
Let $A \in \mathrm{GL}(n, \mathbb{F}_{2})$ be an invertible matrix. Then:

\begin{enumerate}
\item[(a)] $A$ can be reduced to a permutation matrix $P$ by a sequence of
type-3 row operations alone, i.e.\ there exist row operations
$R_1, \ldots, R_k$ such that
\begin{equation}
    \label{eq:basic-elementary-matrix-left-multiplication}
	R_{k} \cdots R_{1} \cdot A = P
	\quad \Longrightarrow \quad
	A = R_{1}^{-1} \cdots R_{k}^{-1} \cdot P.
\end{equation}

\item[(b)] $A$ can be reduced to a permutation matrix $P$ by a sequence of
type-3 row operations $R_1, \ldots, R_k$ together with type-3 column
operations $C_1, \ldots, C_{k'}$, i.e.
\begin{equation}
    \label{eq:basic-celementary-matrix-left-right-multiplication}
	R_{k} \cdots R_{1} \cdot A \cdot C_{1} \cdots C_{k'} = P
	\quad \Longrightarrow \quad
	A = R_{1}^{-1} \cdots R_{k}^{-1} \cdot P \cdot C_{k'}^{-1} \cdots C_{1}^{-1}.
\end{equation}
\end{enumerate}

In both cases, each $R_\ell$ and $C_\ell$ is a type-3 elementary matrix
(adding one row, respectively column, to another), and $P$ is a permutation
matrix realizable by type-1 elementary operations.
\end{theorem}

\begin{proof}
\textbf{(a)} We proceed by Gaussian elimination over $\mathbb{F}_2$, using
row operations only. Since $A$ is invertible, its first column is non-zero,
so some row has a $1$ in the first column; we use this row, together with
type-3 row additions, to clear every other entry in the first column below
it. Repeating this column by column reduces $A$ to an upper-triangular
matrix with $1$'s on the diagonal using only type-3 row operations; a
symmetric elimination step, again using only row operations, then clears
all entries above the diagonal, reducing $A$ to the identity (or permutation in our case). Since some
pivot columns may require selecting a non-zero entry from a lower row, this
process in general also uses row swaps, i.e.\ type-1 operations, which we
collect into a permutation matrix $P$. Writing the composition of all row
operations used as $R_k \cdots R_1$, we obtain $R_k \cdots R_1 \cdot A = P$,
and since every type-3 elementary matrix over $\mathbb{F}_2$ is its own
inverse (as $E(i, j)^2 = I$), inverting each factor gives
$A = R_1^{-1} \cdots R_k^{-1} \cdot P$.

\textbf{(b)} The two-sided case follows by allowing column operations to
participate in the elimination as well. Using row operations, we first
clear all entries below the diagonal, exactly as in part (a); the remaining
above-diagonal entries can then be cleared using column operations instead
of row operations, by adding an already-cleared column into a later one.
Interleaving these two operation types, together with the row and column
swaps needed to fix pivot positions (collected into a single permutation
matrix $P$), reduces $A$ to $P$ via
$R_{k} \cdots R_{1} \cdot A \cdot C_{1} \cdots C_{k'} = P$. Inverting the row
and column factors on each side, using the same self-inverse property as in
part (a), yields
$A = R_{1}^{-1} \cdots R_{k}^{-1} \cdot P \cdot C_{k'}^{-1} \cdots C_{1}^{-1}$.
\end{proof}

\begin{theorem}
\label{thm:elementary-matrix-operation-type-1-and-3}
Let $A \in \mathrm{GL}(n, \mathbb{F}_2)$ be an invertible matrix. Then $A$
can be reduced to a permutation matrix $P$ by a finite sequence of type-3
elementary matrices $E_{3}(i, j)$, applied either from the left
(row operations) or from the right (column operations), together with a
single permutation matrix $P$ realizable by type-1 elementary operations:
\begin{equation}
	\label{eq:elementary-matrix-left-multiplication}
	E_{3}(i_{k}, j_{k}) \cdots E_{3}(i_{1}, j_{1}) \cdot A = P
	\;\Longrightarrow\;
	A = E_{3}(i_{1}, j_{1})^{-1} \cdots E_{3}(i_{k}, j_{k})^{-1} \cdot P
\end{equation}
\begin{align}
	\label{eq:elementary-matrix-left-right-multiplication}
	E_{3}(i_{k}, j_{k}) \cdots E_{3}(i_{1}, j_{1}) \cdot A \cdot E_{3}(i'_{1}, j'_{1}) \cdots E_{3}(i'_{k'}, j'_{k'}) = P
	\\\;\Longrightarrow\;
	A = E_{3}(i_{1}, j_{1})^{-1} \cdots E_{3}(i_{k}, j_{k})^{-1} \cdot P \cdot E_{3}(i'_{k'}, j'_{k'})^{-1} \cdots E_{3}(i'_{1}, j'_{1})^{-1}
\end{align}
where the left-hand factors correspond to row operations and the
right-hand factors correspond to column operations.
\end{theorem}

A CNOT circuit is a quantum circuit composed entirely of CNOT gates. 
Since each CNOT gate performs a linear operation on two qubits, 
the action of an $n$-qubit CNOT circuit as a whole can be 
represented by an invertible $n \times n$ Boolean matrix $A \in GL(n, \mathbb{F}_{2})$. 
Specifically, a CNOT gate controlled by qubit $i$ and targeting qubit 
$j$ corresponds to a type-3 elementary matrix $E_{3}(i, j)$, obtained by adding the $i$-th row to the $j$-th row of the identity matrix $I_{n}$. Applying this gate to a circuit is therefore equivalent to left-multiplying the operator matrix by $E_{3}(i, j)$, i.e., performing the elementary row operation $r_{j} \leftarrow r_{i} \oplus r_{j}$.

In addition to type-3 matrices, we also make use of type-1 elementary matrices $E_{1}(i, j)$, obtained by exchanging the $i$-th and $j$-th rows of $I_{n}$. These correspond to a permutation of two qubits, which can be implemented at zero gate cost via simple rewiring rather than an actual quantum gate.

Since qubit permutations are cost-free under full qubit connectivity, this decomposition directly yields a CNOT circuit implementing $A$: each type-3 factor $E_{3}(i, j)$ corresponds to one CNOT gate, and the total number of factors $k$ determines the CNOT count of the resulting circuit. Crucially, if two elementary operations act on disjoint pairs of qubits, the corresponding CNOT gates can be executed in parallel, meaning the depth of the synthesized circuit depends not on $k$ alone, but on how these operations can be grouped into parallel layers.

This equivalence reframes the problem of CNOT circuit synthesis as a purely linear-algebraic one: transforming an invertible matrix  $A$ into the permutation matrix (or an identity matrix) using the fewest possible sequential layers of row and column operations. All of the synthesis methods reviewed below: Gaussian elimination, DaCSynth, and greedy cost-minimization algorithms, can be understood as different strategies for solving this matrix reduction problem while minimizing its resulting circuit depth.

\section{Schaeffer and Perkowski's Method}
Schaeffer and Perkowski~\cite{schaeffer2014cost} proposed two bidirectional
heuristic methods for synthesizing linear reversible circuits: the
Alternating Elimination with Cost Minimization method (AECM), which extends
Gaussian elimination by diagonalizing rows and columns in a data-dependent
order, and the Multiple CNOT Gate method (MCG), a purely cost-driven search
that evaluates pairs of CNOT gates at each step and is not restricted to the
Gaussian elimination framework. Unlike Patel et al.'s ``Algorithm 1''
\cite{markov2008optimal}, which processes rows and columns in a fixed,
sequential order, both AECM and MCG select row and column operations
adaptively, choosing whichever operation yields the greatest reduction in
cost at each step. Their experiments showed that MCG achieves the lowest
average CNOT gate counts for circuits of up to 24 lines, while AECM performs
best between 28 and 60 lines, and both methods outperform Patel et al.'s
algorithm across most tested circuit sizes.

Central to both methods is a heuristic cost function that measures how far
a matrix $M$ is from the identity matrix $I$, taking into account both $M$
and its inverse:
\begin{equation}
	h_{origin}(M) = \sum_{i=1}^{n}\sum_{j=1}^{n}(M_{(i,j)} \oplus I_{(i,j)}) + (M^{-1}_{(i,j)} \oplus I_{(i,j)})
\end{equation}
The first term counts the number of entries in which $M$ differs from the
identity matrix, while the second term counts the corresponding differences
in $M^{-1}$. Including the inverse term reflects the fact that a row
operation applied to $M$ corresponds to a column operation on $M^{-1}$;
accounting for both therefore gives a more balanced estimate of the number
of elementary operations still required to reduce $M$ to the identity, and
helps guide the search away from states that appear close to complete in
$M$ alone but remain far from complete once the inverse is also considered.
Both AECM and MCG use this cost function to guide their respective row- and
column-selection strategies, committing at each iteration to the operation,
or pair of operations, that most reduces $h_{origin}$.

\section{De Brugière et al.'s Method}
De Brugière et al. \cite{9475957} addressed the problem of minimizing the depth of linear reversible quantum circuits, motivated by the observation that circuit depth is closely tied to execution time under limited qubit decoherence. Their work reformulates the synthesis of a CNOT circuit as the problem of reducing an invertible Boolean matrix to the identity through elementary row and column operations, and proposes two complementary strategies to solve it.

The first, DaCSynth, is a divide-and-conquer framework that recursively partitions the matrix into blocks and eliminates them using a bipartite-graph matching approach, yielding a provable worst-case depth upper bound of $\frac{4}{3}n + 8\lceil \log_2(n) \rceil$ for an $n$-qubit operator, an improvement over prior asymptotic and practical bounds for intermediate-sized registers. The second is a family of purely greedy, cost-minimization algorithms that iteratively apply row and column operations to reduce a chosen cost function toward its minimum, reached when the matrix becomes a permutation matrix. Four cost functions are considered, based on the sum and logarithmic Hamming weight of matrix rows (and their inverses), which balance encouraging overall sparsification against prioritizing ``nearly complete`` rows to avoid local minima in Table \ref{tab:cost-function}.

Both cost functions are trying to answer the same question at each greedy step: "of all the candidate row/column operations available right now, which one should we apply?" Each candidate operation typically reduces some row's Hamming weight by exactly one (e.g., XORing away a 1-bit). The cost function determines how much ``credit`` the algorithm assigns to that reduction and that credit is what steers the greedy search. Consider a row with current Hamming weight $\omega$. Reducing it to $\omega−1$ changes its contribution to the cost function by (1) Under linear cost ($h_{sum}$): the drop is always exactly 1, 
no matter what $\omega$ was. Every row looks equally "worth fixing." (2) Under logarithmic cost ($h_{prod}$, i.e., $log_{2}(\omega)$): the drop is $log_{2}(\omega)-log_{2}(\omega - 1)$, which gets smaller as $\omega$ grows and larger as $\omega$
 shrinks toward 1. This is precisely why $h_{prod}$  gives priority to ``almost done`` rows so that the overall efficiency of elimination is guaranteed, a row that's already nearly at weight 1 gives a bigger log-scale payoff for finishing it off, so the search tends to polish off easy rows first.

\begin{table}
    \centering
    \begin{tabularx}{\textwidth}{C|C}
        \hline
        Cost Function  & Linear Reversible Cost Function\\
         \hline
         $$h_{sum}(A) = \sum_{i,j}a_{i,j}$$ & $$H_{sum}(A) = h_{sum}(A) + h_{sum}(A^{-1})$$\\
         \hline
         $$h_{prod}(A) = \sum_{i=1}^{n}log_{2}(\sum_{j=1}^{n}a_{i,j})$$ & $$H_{prod}(A) = h_{prod}(A) + h_{prod}(A^{-1})$$\\
    \hline
\end{tabularx}
    \caption{Cost function proposed by \cite{10.1145/3474226}}
    \label{tab:cost-function}
\end{table}

Their benchmarks show that while DaCSynth performs consistently well across problem sizes and provides strong worst-case guarantees, the greedy methods outperform DaCSynth for small or structurally simple operators, though their performance degrades sharply as matrix size grows due to the difficulty of escaping local minima in the cost landscape. Applied to a library of reversible functions (including Galois field multipliers and arithmetic circuits), their combined framework achieves depth reductions of up to 92\% (ancilla-free) and 99\% (with ancillary qubits) relative to prior state-of-the-art methods.

This greedy cost-minimization framework forms the direct algorithmic foundation for the depth-oriented approach adopted in this thesis: subsequent work by Shi and Feng \cite{cryptoeprint:2024/381}, and our own enhanced algorithm, build on and extend the cost functions and search strategy first introduced here.

\section{Shi and Feng's Method}

\begin{algorithm}[htbp]
\caption{Improved greedy algorithm for synthesizing low-depth CNOT circuits}
\label{alg:Shi-and-Feng-AES-24}
\footnotesize
\begin{algorithmic}[1]
\Require An invertible matrix $A_{n \times n}$
\Ensure A depth $d$ and a vector of layers $Layers$ that implement $A$ with length $d$
	\State $Layers \leftarrow \varnothing$, $Layers_{r} \leftarrow \varnothing$, $Layers_{c} \leftarrow \varnothing$;
	\State $L_{r} \leftarrow \varnothing$, $L_{c} \leftarrow \varnothing$;
	\State $List \leftarrow \varnothing$;
	\State $B \leftarrow A$;
	\State Randomly determine $H_{r}, H_{c} \leftarrow H_{sqr}, H_{sqc}$ or $H_{prodr}, H_{prodc}$
	\State $cost \leftarrow \{H_{r}(B), H_{c}(B)\}$;
	\State $can\_one \leftarrow$ \textbf{False};
	\State $d \leftarrow 0$;
	\While{True}
		\State $cost \leftarrow\{H_{r}(B), H_{c}(B)\}$;
		\State $mincost \leftarrow$ the minimum resulting cost of all available row operations (can act in parallel with those in $L_{r}$ ) of adding the $i$-th row to the $j$-th row of B(denoted $\{i, j, 0\}$), and if not can one, all available column operations (can act in parallel with those in $L_{c}$ ) adding the $i$-th column to the $j$-th coloumn of $B$ (denoted $\{i, j, 1\}$); $\rhd$ If the current matrix can be implemented with depth 1, only row operations need to be considered.
		\State $List \leftarrow$ {All operations which can lead to the cost of the resulting matrix being mincost};
		\If{$mincost$ == $cost$}
			\If{not $can\_one$ and Can-depth-one($A$)}
				\State $can\_one \leftarrow True$
            \EndIf
            \If{$L_{r}$.size()}
	            \State $d \leftarrow d+1, Layers_{r}$.push\_back($L_{r}$), $L_{r}$.clear();
            \EndIf
            \If{$L_{c}.size()$}
	            \State $d \leftarrow d+1, Layers_{c}$.push\_back($L_{c}$), $L_{c}$.clear();
            \EndIf
            \If{$cost = 2n$}
	            \Break
            \EndIf
            \If{$d >= 100$}
	            \Return $d, \varnothing$
            \EndIf
        \Else
	        \State Randomly choose one operation $\{i, j, op\}$ that minimizes the cost function of the resulting matrix, add $\{i, j, op\}$ to $L_{r}$ if $op == 0$, or $L_{c}$ if $op == 1$;
	        \State $B \leftarrow E(j+i)^{1-op}BE(i+j)^{op}$;
        \EndIf
        \State $List$.clear();
    \EndWhile
    \State Record a permutation $P$ , satisfying $P(i) == j$ if $B[i][j] == 1$;
    \For{$i$ from 0 to $(Layer_{c}.size()-1)$}
	    \State $l \leftarrow \varnothing$;
	    \For{$j$ from 0 to $(Layer_{c}[i].size()-1)$}
		    \State $\{t, c, op\} \leftarrow Layer_{c}[i][j], l$.push\_back($\{c, t\}$);
        \EndFor
        \State $Layers$.push\_back($l$);
    \EndFor
    \For{$i$ from $(Layer_{r}.size()-1)$ down to 0}
	    \State $l \leftarrow \varnothing$;
	    \For{$j$ from 0 to $(Layer_{r}[i].size()-1)$}
		    \State $\{c, t, op\} \leftarrow Layer_{c}[i][j], l$.push\_back($\{P[c], P[t]\}$);
        \EndFor
    \EndFor
\State  \Return $d, Layers$;
\end{algorithmic}
\end{algorithm}

\begin{table}[htbp]
\centering
\begin{tabular}{c|c|c}
\hline
Cost Function & Direction  & Linear Reversible Cost Function\\
\hline
\multirow{2}{*}{$h_{sq}(A) = \sum_{i}(\sum_{j}a_{i,j})^{2}$} & Row & $H_{sqr}(A) = h_{sq}(A) + h_{sq}((A^{-1})^{T})$ \\ 
\cline{2-3}
& Column & $H_{sqc}(A) = h_{sq}(A^{T}) + h_{sq}(A^{-1})$ \\ 
\hline
\end{tabular}
	\caption{Square the Hamming weight of each row to amplify the contribution of incomplete rows proposed by \cite{cryptoeprint:2024/381}}
	\label{tab:square-hamming-weight}
\end{table}

Shi and Feng \cite{cryptoeprint:2024/381} extended the depth-oriented greedy framework of de Brugière et al. to specifically target the synthesis of low-depth CNOT circuits for symmetric-cipher linear layers, with a particular focus on AES MixColumns. Their central algorithmic contribution refines the original greedy cost-minimization strategy in three respects. 

First, in addition to the logarithmic Hamming-weight cost function, they
introduce a cost function based on the square of each row's Hamming weight
(Table~\ref{tab:square-hamming-weight}), which more aggressively prioritizes
rows that are ``far from complete``. As with the previous cost function,
consider a row with current Hamming weight $\omega$: under the squared cost
($h_{sq}$), the drop is $\omega^{2} - (\omega-1)^{2} = 2\omega - 1$, which
increases with $\omega$. A heavy row (say, weight $10$) yields a payoff of
$19$ for a single reduction, while a nearly-done row (weight $2$) yields only
$3$. Thus, the squared cost function does the opposite of the logarithmic
one: it rewards the greedy search far more for chipping away at whichever row
is currently the worst offender.

Their work motivates this with a simple structural observation: the larger the Hamming weight of a row, the more potential gates need to be done on that row, so prioritizing rows or columns with larger Hamming weights may be preferable for lower circuit depth. If a few rows are left ``far from complete`` while everything else is nearly done, those straggler rows become a bottleneck, every remaining layer of the circuit has to wait for them, since depth is governed by whichever row takes the most sequential operations to finish. By squaring the weight, the algorithm is nudged to knock out the heaviest, most-unfinished rows early and in parallel, rather than letting them lag behind while it polishes easier rows which is exactly the failure mode that inflates depth.

So the core idea of Shi and Feng's Algorithm~\ref{alg:Shi-and-Feng-AES-24} is: squaring reshapes the incentive structure of the cost function so that the search actively hunts down the worst-off rows first, rather than treating all rows as equally urgent (linear) or preferring to finish the easy ones first (logarithmic). This is a complementary heuristic to $h_{prod}$'s ``finish the easy ones`` bias, used together (as the paper does, randomly alternating between $H_{sq}$ and $H_{prod}$ variants), the two give the greedy search two different strategies to try, since neither is uniformly the best choice for all matrices.

Second, row and column operations are evaluated with distinct cost functions, 
so that column-wise sparsity is explicitly accounted for during 
column operations rather than only through the logic used for 
rows; finally, they derive a necessary and sufficient condition 
for determining whether a matrix can be implemented with depth 
exactly 1, allowing sparse matrices to be handled more efficiently 
in the final stage of the search in Table \ref{tab:product-hamming-weight}.

\begin{table}[htbp]
\centering
\begin{tabular}{c|c}
\hline
Direction  & Linear Reversible Cost Function\\
\hline
Row & $H_{prodr}(A) = h_{prod}(A) + h_{prod}((A^{-1})^{T})$ \\ 
\hline
Column & $H_{prodc}(A) = h_{prod}(A^{T}) + h_{prod}(A^{-1})$ \\ 
\hline
\end{tabular}
	\caption{Guide matrices that are nearly complete toward the permutation matrix.}
	\label{tab:product-hamming-weight}
\end{table}

Applying this improved algorithm to AES MixColumns, they report a CNOT circuit of depth 10, improving on the previous best result of depth 16 reported by Liu et al.\cite{cryptoeprint:2023/1417}, and also outperforming a direct application of de Brugière et al.'s original greedy algorithm, which yields depth 12 on the same matrix. Beyond AES, they apply their method to a broad collection of MDS matrices and linear layers used in other block ciphers (including ANUBIS\cite{barreto2000anubis}, CLEFIA\cite{shirai2007128}, TWOFISH\cite{schneier1998twofish}, and several lightweight ciphers), consistently achieving the lowest reported depths across nearly all cases. In addition to circuit-depth optimization, the paper proposes a compressed pipeline structure for synthesizing full AES encryption circuits and Grover/Encryption oracles, which reduces the number of intermediate qubit registers relative to the standard pipeline structure while retaining comparable round depth, and reports the lowest published T-depth for an AES-128 Encryption oracle to date.

\section{Clarify Shi and Feng's Method}
In this section, we clarify Shi and Feng's work by presenting it 
as pseudocode in Algorithm~\ref{alg:clarify-Shi-Feng-Greedy}. We 
modify their strategy so that operators which do not reduce the 
cost are not selected initially; only once the stuck threshold 
is exceeded do we accept operations that increase the cost 
function. Additionally, in their "can\_one" condition, they 
consider only column operators for collection; based on our 
observations, we instead change the cost function in this 
condition to Table~\ref{tab:product-hamming-weight}, which 
guides $M$ when it is nearly complete, toward the correct 
direction more quickly. The effect of this modification can be 
seen in our experimental results (Chapter~\ref{ch:experimental-results}). 
And following we detailed introduce our modification, some of function 
will using in further our greedy algorithms.

Since our approach explores all possible operations until $M$ becomes a 
permutation matrix, we adopt a strategy that proceeds from a 
maximum value toward a minimum; accordingly, we initialize our 
comparison condition using system.float\_info.max, called 
``minm``. And we choose The selection list stores the set of 
candidate operations available for the current optimization 
step which called ``select\_list``, visited list tracks the 
selected row and column indices; once a position is chosen, 
it is excluded from future iterations to prevent redundant 
operations. for achieve parallel execution the operations in 
quantum circuit, we have record the used position of CNOT, 
in this case we create two visited lists called $visi_{r}$ and $visi_{c}$.

A termination flag monitors whether the matrix is nearing a 
permutation state. If the current matrix can be transformed 
into a permutation matrix within a single layer, the flag is 
set to True; otherwise, it remains False, it called ``close\_permu``. 
Then we using ``LATEST\_MINM\_DEPTH`` for record the latest minimum depth, if the 
depth over this size, then we setting ``over`` to True.

During the decomposition procedure, we create two large table for 
collection available operations that is $L_{r}$ and $L_{c}$, 
and their cost we estimates then put into $L_{r}^{cost}$ and 
$L_{c}^{cost}$ respectively. To estimates the cost for current 
matrix $M$ that we using ``functionCost()``, not randomly choose
cost function, we use $p$ value to decide which cost function use 
as Algorithm~\ref{alg:functionCost}. This function separate to 
direction and non-direction cost function for matrix $M$ and its
inverse $M'$.

Since the operations on $M$ and $M^{-1}$ are dual, a row operation
applied to $M$ is equivalent to a corresponding column operation applied to
$M^{-1}$, and vice versa, the row-based cost function $h_{sq}$
($h_{prod}$), which measures the squared (logarithmic) Hamming weight of
each row, cannot be applied uniformly to both matrices. When
evaluating the cost of a row operation, the change to $M$ is already
row-wise and can be measured directly by $h_{sq}(M)$ ($h_{prod}(M)$), 
whereas the corresponding change to $M^{-1}$ is column-wise and
must first be converted to a row-wise view by transposition, giving
$h_{sq}((M^{-1})^{\top})$ ($h_{prod}((M^{-1})^{\top})$):
\begin{align}
	H_{r}(M) &= h_{sq}(M) + h_{sq}\!\left((M^{-1})^{\top}\right).\\
             &= h_{prod}(M) + h_{prod}\!\left((M^{-1})^{\top}\right).
\end{align}
Symmetrically, when evaluating the cost of a column operation, the change to
$M$ is column-wise and must be transposed before applying $h_{sq}$ ($h_{prod}$), while
the corresponding change to $M^{-1}$ is already row-wise:
\begin{align}
	H_{c}(M) &= h_{sq}(M^{\top}) + h_{sq}(M^{-1}).\\
             &= h_{prod}(M^{\top}) + h_{prod}(M^{-1}).
\end{align}

Both equations are using ``costMat()`` function to calculate the 
distinct cost value, using $p$ as condition flag that choose cost 
function following all we mentioned approaches, including De 
Brugière et al.'s sum and prod methods, Shi and Feng's square strategy, 
in this work, we also reference the Schaeffer and Perkowski's $h_{origin}$, 
all of them can check on Algorithm~\ref{alg:costMat}.

Since Shi and Feng they are consider the ``can\_one`` to trigger 
Hamming weight of row to Hamming weight to column, our trigger is 
``close\_permu`` and we check whether $M$ has become a permutation 
matrix; since our goal is to find the shortest path to a permutation 
(identity) matrix, we check whether the Hamming weight of each row 
and column of $M$ is less than or equal to $2$. When flag is false, 
we collection row and column simultaneously, otherwise, when flag is 
true, we only collection column operations.

To collect the set of operations available for circuit synthesis, 
we follow the principle that an operator cannot be selected if its 
row or column has already been visited, i.e.\ if it appears in 
$visi_{r}$ or $visi_{c}$. This constraint ensures efficient 
use of the quantum circuit area: for each layer, we aim to place 
as many operators as possible. Our tables $L_r$ and $L_c$ are constructed 
to guarantee this property, as described in Algorithm~\ref{alg:L-collection}.

After collecting all valid operators into our table $L$, we 
select a subset of operations to constitute a single-layer circuit. 
The core idea is that every selected operation should reduce the 
value of the cost function when applied, as shown in Algorithm~\ref{alg:ops-sel}. 
We use the function ``row\_i2j()`` to execute the type-1 elementary row 
(or column) operation of adding one row (or column) to another, 
and then place the result into the target matrix. Note that 
``row\_i2j()`` is applied to $M$ and ``col\_i2j()`` to $M'$ when the operator 
originates from a row; the roles are exchanged when the operator 
originates from a column.

The cost value of every available operator is stored in 
$L_{r}^{cost}$ or $L_{c}^{cost}$ compute with ``operationCost()``, 
which is same with ``functionCost()`` feature depending on the operator 
type. An operator is selected, added to "select\_list" to form 
part of the current layer, and the minimum cost is updated accordingly, 
if its cost value is lower than the current minimum. We also accept 
operators whose cost leaves the cost function unchanged 
(i.e.\ equal to the current value), but in this case the 
minimum cost is not updated.

Occasionally, the search falls into a local minimum. To address 
this, we adopt Shi and Feng's strategy which no longer requires 
the cost to decrease. Instead, we allow the cost to increase temporarily 
when stuck iteration over threshold, which may help the search 
escape the local minimum.

For the operations obtained from ``select\_list``, we execute them on 
$M$ and $M'$ in Algorithm~\ref{alg:ops-exec}, randomly choosing an 
operator that reduces the cost value; this choice is also recorded 
in $visi_{r}$ or $visi_{c}$ to guarantee that the corresponding 
circuit position is not used again.

Another condition arises when ``select\_list`` is empty, meaning that
all positions in the current layer have been filled. In this case, we
transfer $L_r$ into $Ls_r$ and $L_c$ into $Ls_c$ to finalize the synthesis
for this layer, and then reset the visited lists before proceeding to the
next round.

Finally, if the depth of $M$ exceeds the previous minimum depth
``LIMIT``, we set ``over`` to \textbf{True} and break out of the
iteration; otherwise, we return all relevant information for the next
check. Overall, we have clarified Shi and Feng's improved greedy algorithm
by pointing out its limitations, while also incorporating some of the
methods they employ. In the next chapter, we present our proposed
approach, developed in similar spirit.

\begin{algorithm}[htbp]
\caption{Clarified Row/Col Greedy}
\label{alg:clarify-Shi-Feng-Greedy}
\footnotesize
\begin{algorithmic}[1]
\Require $M \in \mathbb{F}_{2}^{n \times n}$, $M' \in \mathbb{F}_{2}^{n \times n}$, $L_{r}, L_{c}, Ls_{r}, Ls_{c}, p$ and over
\Ensure Collection $Ls_{r}, Ls_{c}, OPs_{r}, OPs_{c}$ that implement permutation matrix with length $d$
\State minm $\leftarrow$ Max \Comment{Assign the maximum representable floating-point value as the initial upper-bound.}
\State $visi_{r}, visi_{c}$, select\_list $\leftarrow$ list[$\varnothing$] \Comment{Initialize all required lists.}
\State close\_permu $\leftarrow$ False \Comment{Define the collection condition.}
\State LIMIT $\leftarrow$ LATEST\_MINM\_DEPTH \Comment{Define the previous minimum cost.}
\While{$M$ $\neq$ Permutation} \Comment{Stop when $M$ becomes a permutation matrix.}
	\State select\_list.clear() \Comment{Clear ``select\_list`` for collection in the next iteration.}
	\State $L_{r}$.clear(), $L_{c}$.clear(), $L_{r}^{cost}$.clear(), $L_{c}^{cost}$.clear() \Comment{Clear all lists for the next iteration.}
	
	\State $H_{r}, H_{c}$, minm $\leftarrow$ functionCost($M, M', p$) \Comment{Calculate the current cost of $M$ and $M'$.}
    
    \If{!close\_permu}
		\State $L_{r} \leftarrow L\_Collection(visi_{r})$ \Comment{Collect row operations when $M$ is not close to a low-Hamming-weight state.}
	    \State select\_list, minm $\leftarrow Ops\_Sel( L_{r}, L_{r}^{cost}, M, M', p)$ \Comment{Select operations when $M$ is not close to a permutation matrix.}
    \EndIf

	\State $L_{c} \leftarrow L\_Collection(visi_{c}$) \Comment{Collect column operations; if $M$ is close to a permutation matrix, consider only these operators.}
	\State select\_list, minm $\leftarrow$ Ops\_Sel($L_{c}, L_{c}^{cost}, M, M', p$) \Comment{Selection column operations.}
    
	\State $L_{r}, L_{c}, Ls_{r}, Ls_{c}, M, M', OPs_{r}, OPs_{c}, visi_{r}, visi_{c}, d$, over $\leftarrow$ Ops\_Exec() \Comment{Execute the randomly selected row and column operations from
``select\_list`` on $M$ and $M'$; the selected operators are recorded
in $L_{r}, L_{c}, OPs_{r}, OPs_{c}$, and $visi_{r}, visi_{c}$ are updated
accordingly. Afterward, we check whether $M$ is close to a permutation
matrix; if so, ``close\_permu`` is set to \textbf{True}, and otherwise it remains unchanged.}
    
    \If{$d >$  LIMIT}
	    \State over $\leftarrow True$ \Comment{Break out of the iteration when the depth of $M$ exceeds the previous minimum cost value.}
	    \State \Return
    \EndIf
\EndWhile
\State \Return $L_{r}, L_{c}, Ls_{r}, Ls_{c}, M, OPs_{r}, OPs_{c}, d$, over
\end{algorithmic}
\end{algorithm}

\begin{algorithm}[htbp]
\caption{functionCost}
\label{alg:functionCost}
\begin{algorithmic}[1]
\Require $M \in \mathbb{F}_{2}^{n \times n}, M' \in \mathbb{F}_{2}^{n \times n}, p$
\Ensure Calculate the minimum cost for current $M$ and $M'$ depends on $p$ value
\If{$p ==$ \emph{square} or \emph{prod}} \Comment{Both the square-based and product-based cost functions are considered.}
	\State $H_{r}$ = costMat($M, p$) + costMat($M'^{\top}, p$) \Comment{Calculate the cost for the row direction.}
	\State $H_{c}$ = costMat($M^{\top}, p$) + costMat($M', p$) \Comment{Calculate the cost for the column direction.}
	\State cost $\leftarrow$ max($H_{r}, H_{c}$) \Comment{Choose the maximum value as the upper bound for $M$.}
\Else
	\State cost $\leftarrow$ costMat($M, p$) + costMat($M', p$) \Comment{The sum-based and origin-based cost functions do not consider direction.}
\EndIf
\State \Return cost
\end{algorithmic}
\end{algorithm}

\begin{algorithm}[htbp]
\caption{costMat}
\label{alg:costMat}
\begin{algorithmic}[1]
\Require $M \in \mathbb{F}_{2}^{n \times n}, p$
\Ensure Calculate the cost value for each cost function depends on $p$ value
\State cost $\leftarrow \varnothing$ \Comment{Clear the cost value.}
\State $n \leftarrow |M|$ \Comment{Calculate the length of $M$.}
\Match{$p$}
	\Case{prod}
		\State cost $\leftarrow \sum_{i=0}^{n-1} log_{2}(\sum_{j=0}^{n-1} M_{(i, j)})$ \Comment{Compute the logarithm of each row's elements in $M$.}
	\EndCase
	\Case{sum}
		\State cost $\leftarrow \sum_{i=0}^{n-1}(\sum_{j=0}^{n-1}M_{(i, j)}$) \Comment{Sum all elements in $M$.}
	\EndCase
	\Case{square}
		\State cost $\leftarrow \sum_{i=0}^{n-1}(\sum_{j=0}^{n-1}M_{(i,j)})^{2}$ \Comment{Square the sum of each row's elements in $M$.}
	\EndCase
	\Case{origin}
		\State cost $\leftarrow \sum_{i=0}^{n-1}\sum_{j=0}^{n-1}(M_{(i,j)} \oplus I_{n \times n})$ \Comment{XOR the row elements of $M$ with the identity matrix.}
	\EndCase
\EndMatch
\State \Return cost
\end{algorithmic}
\end{algorithm}

\begin{algorithm}[htbp]
\caption{L\_Collection}
\label{alg:L-collection}
\begin{algorithmic}[1]
\Require $visi_{r}$
\Ensure Capture all of available operators in large table $L$ depends on row operator
\State $n \leftarrow |visi_{r}|$ \Comment{Calculate the length of $visi_{r}$.}
	\For{$i = 0$ \textbf{to} $n-1$}
		\If{$visi_{r}[i] == 1$}
			\State \textbf{continue} \Comment{Skip this control position if its value is 1.}
		\For{$j = 0$ \textbf{to} $n-1$}
			\If{$visi_{r}[j] == 0$ \textbf{and} $j \neq i$}
                \State $L$.append($i, j$) \Comment{Select the operation only if both the control and target positions are unused, and the control and target positions are not the same.}
			\EndIf
		\EndFor
		\EndIf
	\EndFor
\State \Return $L$
\end{algorithmic}
\end{algorithm}

\begin{algorithm}[htbp]
\caption{Ops\_Sel}
\label{alg:ops-sel}
\begin{algorithmic}[1]
\Require $L_{r}, L_{r}^{cost}, M \in \mathbb{F}_{2}^{n \times n}, M' \in \mathbb{F}_{2}^{n \times n}, p$
\Ensure Select all of available operators for a single layer depends on row
\State $n \leftarrow |L_{r}|$ \Comment{Calculate the length of $L_{r}$.}
\State minm $\leftarrow$ MAX \Comment{Assign the maximum representable floating-point value as the initial upper-bound.}
\For{$i = 0$ \textbf{to} $n-1$}
	\State tmp\_mat $\leftarrow$ row\_i2j($M, L_{r}[i]$.control, $L_{r}[i]$.target) \Comment{Apply the corresponding CNOT gate to update $M$.}
	\State tmp\_inv $\leftarrow$ col\_i2j($M', L_{r}[i]$.target, $L_{r}[i]$.control) \Comment{Apply the corresponding CNOT gate to update $M'$.}
	\State cost $\leftarrow$ operationCost(tmp\_mat, tmp\_inv, $p$) \Comment{Recalculate the cost value for $M$.}
	\State $L_{r}^{cost}$.append(cost) \Comment{Record the cost value of current $M$.}
\EndFor
\For{$i = 0$ \textbf{to} $n-1$}
	\If{$L_{r}^{cost}[i] < $ minm} \Comment{If the cost value of $M$ is less than the previous minimum cost value.}
		\State select\_list.clear()
		\State select\_list.append($L_{r}^{cost}[i]$.control, $L_{r}^{cost}[i]$.target) \Comment{Add this operation to ``select\_list``.}
		\State minm $\leftarrow L_{r}^{cost}$.cost \Comment{Update the minimum cost value.}
	\ElsIf{$L_{r}^{cost}[i] == $ minm}
		\State select\_list.append($L_{r}^{cost}[i]$.control, $L_{r}^{cost}[i]$.target) \Comment{Also append operations whose cost remains unchanged.}
	\EndIf
\EndFor
\State \Return select\_list, minm
\end{algorithmic}
\end{algorithm}

\begin{algorithm}[htbp]
\caption{Ops\_Exec}
\label{alg:ops-exec}
\begin{algorithmic}[1]
\Require select\_list, $M \in \mathbb{F}_{2}^{n \times n}, M' \in \mathbb{F}_{2}^{n \times n}, visi_{r}$
\Ensure Execution available operations from select\_list depends on $visi_{r}$ address
\If{|select\_list| $ == 0$}
	\If{$|L_{r}| > 0$}
		\State $Ls_{r}$.append($L_{r}$) \Comment{Add the previous remaining row operations to $Ls_r$.}
		\State $L_{r}$.clear()
		\State $visi_{r}$.clear()
	\EndIf
	\If{can\_close\_permu($M$)}
		\State close\_permu $\leftarrow True$ \Comment{Set ``close\_permu`` to \textbf{True} when $M$ is close to a permutation matrix.}
	\EndIf
\Else
	\State rand $\leftarrow$ random()
	\State sel\_op $\leftarrow$ select\_list[rand] \Comment{Randomly select an operator from select\_list.}
	\State $M \leftarrow$ row\_i2j($M$, sel\_op.control, sel\_op.target) \Comment{Apply the selected CNOT gate to update $M$.}
	\State $M' \leftarrow$ col\_i2j($M'$, sel\_op.target, sel\_op.control) \Comment{Apply the selected CNOT gate to update $M'$.}
	\State $L_{r}$.append(sel\_op.control, sel\_op.target) \Comment{Record the CNOT gate for the current $M$.}
	\If{$\sum visi_{r} == 0$}
		\State $d$++ \Comment{Increase the depth when $visi_{r}$ has been cleared, since a cleared list indicates that it was full in the previous iteration.}
	\EndIf
	\State $visi_{r}$[sel\_op.control] $\leftarrow 1$ \Comment{Record the control position in $visi_r$.}
	\State $visi_{r}$[sel\_op.target] $\leftarrow 1$ \Comment{Record the target position in $visi_{r}$.}
\EndIf
\State \Return $M, M', L_{r}, Ls_{r}, visi_{r}, d$, close\_permu
\end{algorithmic}
\end{algorithm}